%%%%%%%%%%%%%%%%%%%%%%%%%%%%%%%%%%%%%%%%%%%%%%%%%%%%%%%%%%%%%
%% CATEGORY THEORY — Chapters 3–4
%% Limits and Colimits / Adjunctions
%%%%%%%%%%%%%%%%%%%%%%%%%%%%%%%%%%%%%%%%%%%%%%%%%%%%%%%%%%%%%

%%%%%%%%%%%%%%%%%%%%%%%%%%%%%%%%%%%%%%%%%%%%%%%%%%%%%%%%%%%%
%%           CHAPTER 3: LIMITS AND COLIMITS               %%
%%%%%%%%%%%%%%%%%%%%%%%%%%%%%%%%%%%%%%%%%%%%%%%%%%%%%%%%%%%%

\chapter{Limits and Colimits}\label{ch:limits-colimits}
\minitoc

\vspace{1em}

Limits and colimits are the categorical generalisations of constructions
that pervade all of mathematics: products, intersections, kernels,
pullbacks, inverse limits, and their duals.  Understanding them in full
generality is one of the main rewards of learning category theory.
In this chapter we develop the theory systematically, beginning with the
language of diagrams and cones, proceeding to the universal property that
defines a limit, and then examining the most important special cases.
We prove the fundamental existence theorem---that products and equalisers
suffice for all finite limits---and establish the crucial relationship
between limits and representable functors.

\index{limit|(}
\index{colimit|(}

%% ============================================================
\section{Diagrams and index categories}\label{sec:diagrams}
%% ============================================================

\begin{definition}[Diagram]\label{def:diagram}
\index{diagram}
\index{index category}
\index{functor!diagram}
Let $\mathcal{C}$ be a category and let $\mathcal{J}$ be a small category
(the \emph{index category} or \emph{shape category}).
A \textbf{diagram of shape~$\mathcal{J}$ in~$\mathcal{C}$} is a functor
\[
    F \colon \mathcal{J} \longrightarrow \mathcal{C}.
\]
We write $F_j$ for the image of an object $j \in \Ob(\mathcal{J})$
and $F_\alpha \colon F_j \to F_k$ for the image of a morphism
$\alpha \colon j \to k$ in~$\mathcal{J}$.
\end{definition}

\begin{example}\label{ex:diagrams-shapes}
\index{discrete category}
\index{parallel pair}
\index{span}
\index{cospan}
The following small categories appear constantly as index categories.
\begin{enumerate}[label=(\roman*)]
\item \textbf{Discrete categories.}
    Let $\mathcal{J}$ be a discrete category on a set~$I$
    (objects~$I$, only identity morphisms).
    A diagram $F\colon \mathcal{J}\to\mathcal{C}$ is simply a family
    of objects $\{F_i\}_{i\in I}$.

\item \textbf{The parallel pair.}
    Let $\mathcal{J} = (\bullet \rightrightarrows \bullet)$,
    i.e.\ two objects $0,1$ with two non-identity morphisms
    $\alpha,\beta\colon 0\to 1$.
    A diagram of this shape is a pair of parallel morphisms
    $f,g\colon A\to B$ in~$\mathcal{C}$.

\item \textbf{The span.}
    $\mathcal{J} = (\bullet \leftarrow \bullet \rightarrow \bullet)$.
    A diagram is a span $A \xleftarrow{f} C \xrightarrow{g} B$.

\item \textbf{The cospan.}
    $\mathcal{J} = (\bullet \rightarrow \bullet \leftarrow \bullet)$.
    A diagram is a cospan $A \xrightarrow{f} C \xleftarrow{g} B$.
\end{enumerate}
\end{example}

\begin{remark}\label{rem:diagram-notation}
\index{diagram!as functor}
We shall sometimes write a diagram simply as
$\{F_j,\,F_\alpha\}_{j\in\mathcal{J}}$
when the index category is understood.
A diagram is nothing more than a functor; the terminology ``diagram''
merely signals that we are about to take its limit or colimit.
\end{remark}


%% ============================================================
\section{Cones and cocones}\label{sec:cones}
%% ============================================================

\begin{definition}[Cone]\label{def:cone}
\index{cone}
\index{cone!vertex}
\index{cone!leg}
Let $F\colon \mathcal{J}\to\mathcal{C}$ be a diagram.
A \textbf{cone over~$F$} is a pair $(N,\,\{\pi_j\}_{j\in\mathcal{J}})$
consisting of an object $N\in\mathcal{C}$ (the \emph{vertex} of the cone)
and a family of morphisms $\pi_j\colon N\to F_j$ (the \emph{legs}),
one for each object $j\in\mathcal{J}$,
such that for every morphism $\alpha\colon j\to k$ in~$\mathcal{J}$
the following triangle commutes:
\[
\begin{tikzcd}
& N \arrow[dl, "\pi_j"'] \arrow[dr, "\pi_k"] & \\
F_j \arrow[rr, "F_\alpha"'] & & F_k
\end{tikzcd}
\]
That is, $F_\alpha \circ \pi_j = \pi_k$ for every $\alpha\colon j\to k$.
\end{definition}

\begin{definition}[Morphism of cones]\label{def:cone-morphism}
\index{cone!morphism}
Let $(N,\pi_j)$ and $(N',\pi'_j)$ be two cones over $F$.
A \textbf{morphism of cones} from $(N,\pi_j)$ to $(N',\pi'_j)$
is a morphism $u\colon N\to N'$ in~$\mathcal{C}$ such that
$\pi'_j \circ u = \pi_j$ for every $j\in\mathcal{J}$:
\[
\begin{tikzcd}
N \arrow[rr, "u"] \arrow[dr, "\pi_j"'] && N' \arrow[dl, "\pi'_j"] \\
& F_j &
\end{tikzcd}
\]
The cones over~$F$ and their morphisms form a category,
which we denote $\mathsf{Cone}(F)$.
\index{cone!category of cones}
\end{definition}

\begin{definition}[Cocone]\label{def:cocone}
\index{cocone}
\index{cocone!vertex}
\index{cocone!leg}
Dually, a \textbf{cocone under~$F$} is a pair
$(N,\,\{\iota_j\}_{j\in\mathcal{J}})$ where $\iota_j\colon F_j\to N$
and $\iota_k \circ F_\alpha = \iota_j$ for every $\alpha\colon j\to k$:
\[
\begin{tikzcd}
F_j \arrow[rr, "F_\alpha"] \arrow[dr, "\iota_j"'] & & F_k \arrow[dl, "\iota_k"] \\
& N &
\end{tikzcd}
\]
Cocones under~$F$ form a category $\mathsf{Cocone}(F)$.
\end{definition}

\begin{remark}\label{rem:cone-as-nat}
\index{constant functor}
\index{natural transformation!and cones}
Let $\Delta_N \colon \mathcal{J}\to\mathcal{C}$ denote the constant functor
sending every object to~$N$ and every morphism to~$\id_N$.
Then a cone $(N,\pi_j)$ over~$F$ is exactly a natural transformation
$\pi\colon \Delta_N \Rightarrow F$.
Dually, a cocone is a natural transformation $\iota\colon F\Rightarrow\Delta_N$.
This reformulation will be important when we study the relationship
between limits and the functor category $\Fun(\mathcal{J},\mathcal{C})$.
\end{remark}


%% ============================================================
\section{Limits}\label{sec:limits}
%% ============================================================

\begin{definition}[Limit]\label{def:limit}
\index{limit!definition}
\index{universal cone}
\index{limit!universal property}
Let $F\colon\mathcal{J}\to\mathcal{C}$ be a diagram.
A \textbf{limit} of~$F$ is a cone $(\varprojlim F,\,\{\pi_j\}_{j})$
that is \emph{terminal} in $\mathsf{Cone}(F)$.
Explicitly, this means:
\begin{enumerate}[label=(\alph*)]
\item $(\varprojlim F,\,\pi_j)$ is a cone over~$F$; and
\item for every cone $(N,\,\psi_j)$ over~$F$, there exists a
      \emph{unique} morphism $u\colon N\to\varprojlim F$
      such that $\pi_j\circ u = \psi_j$ for every $j\in\mathcal{J}$:
\end{enumerate}
\[
\begin{tikzcd}
N \arrow[dr, "\psi_j"'] \arrow[rr, dashed, "\exists!\, u"] & & \varprojlim F \arrow[dl, "\pi_j"] \\
& F_j &
\end{tikzcd}
\]
We also write $\lim_{j\in\mathcal{J}} F_j$ or simply $\lim F$.
The morphisms $\pi_j$ are called the \textbf{projection morphisms}
(or \textbf{canonical projections}).
\index{projection morphisms}
\end{definition}

\begin{proposition}[Uniqueness of limits]\label{prop:limit-unique}
\index{limit!uniqueness}
If $F\colon\mathcal{J}\to\mathcal{C}$ admits a limit, then the limit
is unique up to unique isomorphism.
\end{proposition}

\begin{proof}
This is a standard argument for terminal objects.
Suppose $(L,\pi_j)$ and $(L',\pi'_j)$ are both limits of~$F$.
By the universal property of~$L$, there is a unique
$u\colon L'\to L$ with $\pi_j\circ u = \pi'_j$ for all~$j$.
By the universal property of~$L'$, there is a unique
$v\colon L\to L'$ with $\pi'_j\circ v = \pi_j$ for all~$j$.
Then $\pi_j\circ(u\circ v) = \pi'_j\circ v = \pi_j$
for all~$j$, and by uniqueness applied to $L$ with the cone $(L,\pi_j)$,
we get $u\circ v = \id_L$.  Similarly $v\circ u = \id_{L'}$.
Hence $u$ is an isomorphism, and it is the unique isomorphism of cones.
\[
\begin{tikzcd}
L \arrow[rr, bend left, "v"] & & L' \arrow[ll, bend left, "u"] \\
& F_j \arrow[ul, leftarrow, "\pi_j"] \arrow[ur, leftarrow, "\pi'_j"'] &
\end{tikzcd}
\]
\end{proof}


%% ============================================================
\section{Colimits}\label{sec:colimits}
%% ============================================================

\begin{definition}[Colimit]\label{def:colimit}
\index{colimit!definition}
\index{universal cocone}
\index{colimit!universal property}
Let $F\colon\mathcal{J}\to\mathcal{C}$ be a diagram.
A \textbf{colimit} of~$F$ is a cocone $(\varinjlim F,\,\{\iota_j\}_{j})$
that is \emph{initial} in $\mathsf{Cocone}(F)$.
That is, for every cocone $(N,\,\psi_j)$ under~$F$, there exists a
unique morphism $u\colon \varinjlim F\to N$
with $u\circ\iota_j = \psi_j$ for all~$j$:
\[
\begin{tikzcd}
F_j \arrow[rr, "\iota_j"] \arrow[dr, "\psi_j"'] & & \varinjlim F \arrow[dl, dashed, "\exists!\, u"] \\
& N &
\end{tikzcd}
\]
The morphisms $\iota_j$ are the \textbf{coprojection morphisms}
(or \textbf{canonical injections}).
\index{coprojection morphisms}
\end{definition}

\begin{remark}\label{rem:colimit-dual}
\index{duality!limits and colimits}
A colimit of $F\colon\mathcal{J}\to\mathcal{C}$
is the same as a limit of $F^{\op}\colon\mathcal{J}^{\op}\to\mathcal{C}^{\op}$.
Hence every theorem about limits has a dual statement about colimits.
We will exploit this duality systematically.
\end{remark}


%% ============================================================
\section{Products and coproducts}\label{sec:products-coproducts}
%% ============================================================

Products and coproducts are limits and colimits indexed by a
discrete category.

\begin{definition}[Product]\label{def:product}
\index{product}
\index{product!universal property}
Let $\{A_i\}_{i\in I}$ be a family of objects in~$\mathcal{C}$,
viewed as a diagram $F\colon I_{\mathrm{disc}}\to\mathcal{C}$.
The \textbf{product} $\prod_{i\in I} A_i$ is the limit of this diagram.
Concretely, it is an object $\prod_{i} A_i$ equipped with projections
$\pi_i\colon \prod_j A_j\to A_i$ such that for every object~$N$
with morphisms $f_i\colon N\to A_i$, there exists a unique
$\langle f_i \rangle\colon N\to\prod_j A_j$ with
$\pi_i\circ\langle f_i \rangle = f_i$:
\[
\begin{tikzcd}
& N \arrow[dl, "f_1"'] \arrow[d, "f_2"] \arrow[dr, "f_3"]
      \arrow[dd, dashed, "\langle f_i\rangle" description] & \\
A_1 & A_2 & A_3 \\
& \prod_j A_j \arrow[ul, "\pi_1"] \arrow[u, "\pi_2"'] \arrow[ur, "\pi_3"'] &
\end{tikzcd}
\]
For a binary product we write $A\times B$ with projections
$\pi_1\colon A\times B\to A$ and $\pi_2\colon A\times B\to B$.
\index{product!binary}
\end{definition}

\begin{definition}[Coproduct]\label{def:coproduct}
\index{coproduct}
\index{coproduct!universal property}
Dually, the \textbf{coproduct} $\coprod_{i\in I} A_i$ is the
colimit of a family indexed by a discrete category.
It is an object with injections $\iota_i\colon A_i\to\coprod_j A_j$
such that for every object~$N$ with morphisms $g_i\colon A_i\to N$,
there is a unique $[g_i]\colon\coprod_j A_j\to N$
with $[g_i]\circ\iota_i = g_i$:
\[
\begin{tikzcd}
A_1 \arrow[dr, "g_1"] \arrow[d, "\iota_1"'] &
A_2 \arrow[d, "\iota_2"] \arrow[d] &
A_3 \arrow[dl, "g_3"'] \arrow[d, "\iota_3"] \\
\coprod_j A_j \arrow[r, dashed, "{[g_i]}"'] & N &
\end{tikzcd}
\]
For a binary coproduct we write $A\amalg B$ or $A + B$.
\index{coproduct!binary}
\end{definition}

\begin{example}[Products and coproducts in concrete categories]%
\label{ex:products-concrete}
\index{product!in Set@in $\Set$}
\index{product!in Grp@in $\Grp$}
\index{product!in Top@in $\Top$}
\index{product!in Ab@in $\Ab$}
\index{coproduct!in Set@in $\Set$}
\index{coproduct!in Grp@in $\Grp$}
\index{coproduct!in Top@in $\Top$}
\index{coproduct!in Ab@in $\Ab$}
\begin{enumerate}[label=(\roman*)]
\item \textbf{$\Set$.} The product is the Cartesian product
    $\prod_{i} A_i = \{\,(a_i)_{i\in I} : a_i\in A_i\,\}$
    with the usual projections.
    The coproduct is the disjoint union $\coprod_{i} A_i$.

\item \textbf{$\Grp$.} The product is the direct product of groups
    (Cartesian product with component-wise operations).
    The coproduct is the \emph{free product} $\ast_{i} G_i$,
    which is considerably more complicated than the direct product.
    \index{free product}

\item \textbf{$\Ab$.} Both the product and the finite coproduct
    coincide with the direct sum $\bigoplus_{i} A_i$
    (for finite families).  For infinite families the product
    is the full direct product and the coproduct is the direct sum
    (elements with finitely many nonzero components).
    \index{direct sum}

\item \textbf{$\Top$.} The product is the Cartesian product with the
    product topology (the coarsest topology making all projections
    continuous).  The coproduct is the disjoint union with the
    coproduct (disjoint union) topology.
    \index{product topology}

\item \textbf{$\Vect_\K$.} The product is the direct product
    of vector spaces and the finite coproduct is the direct sum.
    For infinite families, $\prod_i V_i \neq \bigoplus_i V_i$
    in general.
\end{enumerate}
\end{example}

\begin{exercise}\label{exer:product-comm-assoc}
\index{product!commutativity}
\index{product!associativity}
Let $\mathcal{C}$ be a category with binary products.
Show that:
\begin{enumerate}[label=(\alph*)]
\item $A \times B \cong B \times A$ (commutativity);
\item $(A \times B) \times C \cong A \times (B \times C)$ (associativity);
\item if $\mathcal{C}$ has a terminal object~$1$, then $A \times 1 \cong A$.
\end{enumerate}
\end{exercise}


%% ============================================================
\section{Equalisers and coequalisers}\label{sec:equalisers}
%% ============================================================

\begin{definition}[Equaliser]\label{def:equaliser}
\index{equaliser}
\index{equaliser!universal property}
Let $f,g\colon A\rightrightarrows B$ be a parallel pair.
The \textbf{equaliser} of $f$ and~$g$ is the limit of the diagram
$(\bullet\rightrightarrows\bullet)$: an object $E$ with a morphism
$e\colon E\to A$ such that $f\circ e = g\circ e$,
and universal with this property.
That is, for every $h\colon N\to A$ with $f\circ h = g\circ h$,
there exists a unique $u\colon N\to E$ with $e\circ u = h$:
\[
\begin{tikzcd}
N \arrow[dr, dashed, "\exists!\,u"'] \arrow[drr, bend left, "h"] & & \\
& E \arrow[r, "e"] & A \arrow[r, shift left, "f"] \arrow[r, shift right, "g"'] & B
\end{tikzcd}
\]
We write $E = \eq(f,g)$.
\end{definition}

\begin{example}\label{ex:equaliser-set}
\index{equaliser!in Set@in $\Set$}
In~$\Set$, the equaliser of $f,g\colon A\rightrightarrows B$ is
$\eq(f,g) = \{a\in A : f(a) = g(a)\}$
with the inclusion $e\colon\eq(f,g)\hookrightarrow A$.
\end{example}

\begin{proposition}\label{prop:equaliser-mono}
\index{equaliser!is monic}
\index{monomorphism!and equalisers}
Every equaliser is a monomorphism.
\end{proposition}

\begin{proof}
Let $e\colon E\to A$ be the equaliser of $f,g$.
Suppose $e\circ u = e\circ v$ for $u,v\colon N\to E$.
Then $f\circ(e\circ u) = f\circ(e\circ v)$ and
$g\circ(e\circ u) = g\circ(e\circ v)$.
Since $f\circ e = g\circ e$, the morphism $e\circ u = e\circ v$
satisfies $f\circ(e\circ u) = g\circ(e\circ u)$.
By the universal property of the equaliser applied to $h = e\circ u$,
there is a \emph{unique} morphism $N\to E$ making the triangle commute;
both $u$ and~$v$ satisfy this, so $u = v$.
\end{proof}

\begin{definition}[Coequaliser]\label{def:coequaliser}
\index{coequaliser}
\index{coequaliser!universal property}
Dually, the \textbf{coequaliser} of $f,g\colon A\rightrightarrows B$
is an object $Q$ with $q\colon B\to Q$ such that $q\circ f = q\circ g$,
universal from below:
\[
\begin{tikzcd}
A \arrow[r, shift left, "f"] \arrow[r, shift right, "g"'] &
B \arrow[r, "q"] \arrow[rrd, bend right, "h"'] &
Q \arrow[rd, dashed, "\exists!\,u"] & \\
& & & N
\end{tikzcd}
\]
We write $Q = \coeq(f,g)$.
\end{definition}

\begin{example}\label{ex:coequaliser-set}
\index{coequaliser!in Set@in $\Set$}
In~$\Set$, the coequaliser of $f,g\colon A\rightrightarrows B$ is the
quotient $B/{\sim}$ where $\sim$ is the equivalence relation generated
by $f(a)\sim g(a)$ for all $a\in A$.
\end{example}

\begin{proposition}\label{prop:coequaliser-epi}
\index{coequaliser!is epic}
\index{epimorphism!and coequalisers}
Every coequaliser is an epimorphism.
\end{proposition}

\begin{proof}
Dual to the proof that equalisers are monic (Proposition~\ref{prop:equaliser-mono}).
\end{proof}


%% ============================================================
\section{Kernels and cokernels}\label{sec:kernels-cokernels}
%% ============================================================

\begin{definition}[Kernel]\label{def:kernel-categorical}
\index{kernel!categorical}
\index{zero morphism}
Let $\mathcal{C}$ be a category with a zero object~$0$.
For a morphism $f\colon A\to B$, the \textbf{kernel} of~$f$ is the
equaliser of~$f$ and the zero morphism $0_{AB}\colon A\to 0\to B$:
\[
\Ker(f) = \eq(f,\,0_{AB}).
\]
Concretely, it is a morphism $k\colon K\to A$ such that $f\circ k = 0$
and every $h\colon N\to A$ with $f\circ h = 0$ factors uniquely through~$k$:
\[
\begin{tikzcd}
N \arrow[dr, dashed, "\exists!\,u"'] \arrow[drr, bend left, "h"] & & \\
& K \arrow[r, "k"] & A \arrow[r, "f"] & B
\end{tikzcd}
\]
\end{definition}

\begin{definition}[Cokernel]\label{def:cokernel-categorical}
\index{cokernel!categorical}
The \textbf{cokernel} of $f\colon A\to B$ is the coequaliser of~$f$
and~$0_{AB}$:
\[
\coker(f) = \coeq(f,\,0_{AB}).
\]
It is a morphism $c\colon B\to Q$ with $c\circ f = 0$, universal
with this property.
\end{definition}

\begin{example}\label{ex:ker-coker-ab}
\index{kernel!in Ab@in $\Ab$}
\index{cokernel!in Ab@in $\Ab$}
In~$\Ab$, the kernel of a homomorphism $f\colon A\to B$ is
$\Ker(f) = \{a\in A : f(a) = 0\}$ with the inclusion into~$A$,
and the cokernel is $\coker(f) = B/\im(f)$ with the projection.
\end{example}

\begin{exercise}\label{exer:kernel-mono}
\index{kernel!is monic}
Show that in any category with a zero object, a kernel is a monomorphism.
(This follows from the fact that equalisers are monic.)
\end{exercise}


%% ============================================================
\section{Pullbacks and pushouts}\label{sec:pullbacks-pushouts}
%% ============================================================

\begin{definition}[Pullback]\label{def:pullback}
\index{pullback}
\index{pullback!universal property}
\index{fibre product}
Given a cospan $A \xrightarrow{f} C \xleftarrow{g} B$, the
\textbf{pullback} (or \textbf{fibre product}) is the limit of this
diagram.  It is an object $A\times_C B$ equipped with morphisms
$p_1\colon A\times_C B\to A$ and $p_2\colon A\times_C B\to B$
such that $f\circ p_1 = g\circ p_2$, and universal with this property:
\[
\begin{tikzcd}
N \arrow[drr, bend left, "\psi_2"]
  \arrow[ddr, bend right, "\psi_1"']
  \arrow[dr, dashed, "\exists!\,u" description] & & \\
& A\times_C B \arrow[r, "p_2"] \arrow[d, "p_1"']
    \arrow[dr, phantom, "\lrcorner" very near start]
  & B \arrow[d, "g"] \\
& A \arrow[r, "f"'] & C
\end{tikzcd}
\]
The symbol $\lrcorner$ in the square indicates that it is a pullback square.
\end{definition}

\begin{example}[Pullback in $\Set$]\label{ex:pullback-set}
\index{pullback!in Set@in $\Set$}
Given $f\colon A\to C$ and $g\colon B\to C$ in $\Set$,
\[
    A\times_C B = \{(a,b)\in A\times B : f(a) = g(b)\}.
\]
The projections are $(a,b)\mapsto a$ and $(a,b)\mapsto b$.
\end{example}

\begin{example}[Pullback in $\Top$]\label{ex:pullback-top}
\index{pullback!in Top@in $\Top$}
In $\Top$, the pullback is the set-theoretic fibre product
$\{(a,b)\in A\times B : f(a) = g(b)\}$ equipped with the
subspace topology inherited from $A\times B$.
\end{example}

\begin{example}[Pullback as inverse image]\label{ex:pullback-preimage}
\index{pullback!as inverse image}
In~$\Set$, consider $f\colon A\to C$ and the inclusion
$\iota\colon S\hookrightarrow C$ of a subset.
The pullback $A\times_C S$ is the preimage
$f^{-1}(S) = \{a\in A : f(a)\in S\}$.
\[
\begin{tikzcd}
f^{-1}(S) \arrow[r, hook] \arrow[d]
  \arrow[dr, phantom, "\lrcorner" very near start]
& A \arrow[d, "f"] \\
S \arrow[r, hook, "\iota"'] & C
\end{tikzcd}
\]
\end{example}

\begin{proposition}\label{prop:pullback-of-mono}
\index{pullback!of monomorphism}
\index{monomorphism!stable under pullback}
If $g\colon B\to C$ is a monomorphism, then $p_1\colon A\times_C B\to A$
is also a monomorphism.  That is, monomorphisms are stable under pullback.
\end{proposition}

\begin{proof}
Suppose $p_1\circ u = p_1\circ v$ for $u,v\colon N\to A\times_C B$.
Then $f\circ p_1\circ u = f\circ p_1\circ v$, hence
$g\circ p_2\circ u = g\circ p_2\circ v$, hence
$p_2\circ u = p_2\circ v$ since $g$ is monic.
Now $u$ and $v$ are both morphisms $N\to A\times_C B$ satisfying
$p_1\circ u = p_1\circ v$ and $p_2\circ u = p_2\circ v$.
By the uniqueness part of the pullback universal property, $u = v$.
\end{proof}

\begin{definition}[Pushout]\label{def:pushout}
\index{pushout}
\index{pushout!universal property}
\index{amalgamated sum}
Given a span $A \xleftarrow{f} C \xrightarrow{g} B$, the
\textbf{pushout} (or \textbf{amalgamated sum}) is the colimit of
this diagram.  It is an object $A\amalg_C B$ with morphisms
$q_1\colon A\to A\amalg_C B$ and $q_2\colon B\to A\amalg_C B$
such that $q_1\circ f = q_2\circ g$, universal with this property:
\[
\begin{tikzcd}
C \arrow[r, "g"] \arrow[d, "f"'] & B \arrow[d, "q_2"]
  \arrow[ddr, bend left, "\psi_2"] & \\
A \arrow[r, "q_1"'] \arrow[drr, bend right, "\psi_1"']
  \arrow[dr, phantom, "\ulcorner" very near end]
& A\amalg_C B \arrow[dr, dashed, "\exists!\,u" description] & \\
& & N
\end{tikzcd}
\]
\end{definition}

\begin{example}[Pushout in $\Set$]\label{ex:pushout-set}
\index{pushout!in Set@in $\Set$}
Given a span $A\xleftarrow{f} C \xrightarrow{g} B$ in~$\Set$,
\[
    A\amalg_C B = (A\amalg B)/{\sim}
\]
where $\sim$ is generated by $f(c)\sim g(c)$ for all $c\in C$.
\end{example}

\begin{example}[Pushout in $\Top$]\label{ex:pushout-top}
\index{pushout!in Top@in $\Top$}
\index{attaching space}
In~$\Top$, pushouts correspond to gluing constructions.
If $f\colon A\hookrightarrow X$ is a subspace inclusion and
$g\colon A\to Y$ is continuous, then $X\amalg_A Y$ is the space
obtained by gluing $Y$ to~$X$ along~$A$ via~$g$.
A fundamental instance: if $D^n$ is the $n$-disc,
$S^{n-1}$ its boundary, and $\varphi\colon S^{n-1}\to X$ an attaching map,
then
\[
\begin{tikzcd}
S^{n-1} \arrow[r, hook] \arrow[d, "\varphi"'] & D^n \arrow[d] \\
X \arrow[r] & X\cup_\varphi D^n
\end{tikzcd}
\]
is a pushout square giving the space $X$ with an $n$-cell attached.
\end{example}

\begin{exercise}\label{exer:pushout-epi}
\index{pushout!of epimorphism}
Show that epimorphisms are stable under pushout.
(This is dual to Proposition~\ref{prop:pullback-of-mono}.)
\end{exercise}

\begin{exercise}\label{exer:pullback-pasting}
\index{pullback!pasting lemma}
\textbf{(Pullback pasting lemma.)}
Consider a diagram
\[
\begin{tikzcd}
A \arrow[r] \arrow[d] & B \arrow[r] \arrow[d] & C \arrow[d] \\
D \arrow[r] & E \arrow[r] & F
\end{tikzcd}
\]
Show that if both squares are pullbacks, then the outer rectangle
is a pullback.  Show also that if the right square and the outer
rectangle are pullbacks, then the left square is a pullback.
\end{exercise}


%% ============================================================
\section{General existence theorem}\label{sec:limit-existence}
%% ============================================================

The following fundamental theorem reduces the existence of all finite
limits to two basic building blocks.

\begin{theorem}[Products and equalisers give all finite limits]%
\label{thm:prod-eq-limits}
\index{limit!existence theorem}
\index{finite limits}
\index{completeness!finite}
A category $\mathcal{C}$ has all finite limits if and only if it has all
finite products and all equalisers.
\end{theorem}

\begin{proof}
The ``only if'' direction is clear, since products and equalisers are
special cases of finite limits.

For the ``if'' direction, let $F\colon\mathcal{J}\to\mathcal{C}$
be a diagram where $\mathcal{J}$ is a finite category.
We construct $\lim F$ as an equaliser of two maps between products.

\medskip\noindent
\textbf{Step 1.} Form the products
\[
    P = \prod_{j\in\Ob(\mathcal{J})} F_j,
    \qquad
    Q = \prod_{\alpha\in\Mor(\mathcal{J})} F_{\mathrm{cod}(\alpha)},
\]
where for each morphism $\alpha\colon j\to k$ in~$\mathcal{J}$,
the factor of~$Q$ indexed by~$\alpha$ is $F_k = F_{\mathrm{cod}(\alpha)}$.

\medskip\noindent
\textbf{Step 2.} Define two morphisms $s,t\colon P\to Q$ as follows.
For each morphism $\alpha\colon j\to k$ in~$\mathcal{J}$, the
$\alpha$-component of~$s$ is the composite
\[
    s_\alpha\colon P \xrightarrow{\pi_j} F_j \xrightarrow{F_\alpha} F_k,
\]
and the $\alpha$-component of~$t$ is
\[
    t_\alpha\colon P \xrightarrow{\pi_k} F_k.
\]

\[
\begin{tikzcd}
\eq(s,t) \arrow[r, "e"] & P \arrow[r, shift left, "s"] \arrow[r, shift right, "t"'] & Q
\end{tikzcd}
\]

\medskip\noindent
\textbf{Step 3.} Let $e\colon E\to P$ be the equaliser of $s$ and~$t$.
We claim $E = \lim F$ with projections
$\pi_j\circ e\colon E\to F_j$.

The condition $s\circ e = t\circ e$ means that for every
$\alpha\colon j\to k$,
\[
    F_\alpha \circ \pi_j \circ e = \pi_k \circ e,
\]
which is exactly the cone condition.

Conversely, if $(N,\psi_j)$ is any cone over~$F$, the universal property
of~$P$ gives a unique $h\colon N\to P$ with $\pi_j\circ h = \psi_j$.
The cone condition ensures $s\circ h = t\circ h$, so the universal
property of the equaliser gives a unique $u\colon N\to E$ with
$e\circ u = h$.
Then $\pi_j\circ e\circ u = \pi_j\circ h = \psi_j$, as required.
\end{proof}

\begin{corollary}\label{cor:complete-category}
\index{complete category}
\index{cocomplete category}
A category $\mathcal{C}$ is \textbf{(finitely) complete} if and only if
it has (finite) products and equalisers.  Dually, $\mathcal{C}$ is
\textbf{(finitely) cocomplete} if and only if it has (finite) coproducts
and coequalisers.
\end{corollary}

\begin{example}\label{ex:complete-categories}
\index{complete category!examples}
The categories $\Set$, $\Grp$, $\Ab$, $\Ring$, $\Top$, $\Vect_\K$,
and $\Mod_R$ (for any ring~$R$) are all complete and cocomplete.
\end{example}

\begin{remark}\label{rem:limits-general}
\index{limit!small}
The same argument works for arbitrary (small) limits: a category
with all small products and equalisers has all small limits.
A category with all small limits is called \textbf{complete}.
\end{remark}


%% ============================================================
\section{Limits and the Hom functor}\label{sec:limits-hom}
%% ============================================================

\begin{theorem}[Limits commute with Hom]\label{thm:lim-hom}
\index{limit!and Hom functor}
\index{Hom functor!preserves limits}
Let $F\colon\mathcal{J}\to\mathcal{C}$ be a diagram admitting a limit
$L = \lim F$.  For any object $X\in\mathcal{C}$,
\[
    \Hom_{\mathcal{C}}(X,\,\lim_{j} F_j)
    \;\cong\;
    \lim_{j}\, \Hom_{\mathcal{C}}(X,\,F_j),
\]
naturally in~$X$.  Here the limit on the right is computed in $\Set$.
\end{theorem}

\begin{proof}
Let $(L,\pi_j)$ be the limit cone.  Define
\[
    \Phi\colon \Hom(X,L) \longrightarrow \lim_j\Hom(X,F_j)
\]
by $\Phi(h) = (\pi_j\circ h)_{j\in\mathcal{J}}$.  We verify that
$(\pi_j\circ h)_j$ lies in the limit: for $\alpha\colon j\to k$,
$F_\alpha\circ\pi_j\circ h = \pi_k\circ h$ since $(L,\pi_j)$ is a cone.

The map $\Phi$ is \textbf{injective}: if $\Phi(h)=\Phi(h')$,
then $\pi_j\circ h = \pi_j\circ h'$ for all~$j$, so $h = h'$
by the uniqueness in the universal property of~$L$.

The map $\Phi$ is \textbf{surjective}: given
$(f_j)_j\in\lim_j\Hom(X,F_j)$, the family $(f_j\colon X\to F_j)_j$
is a cone over~$F$ with vertex~$X$.
By the universal property of~$L$, there exists a unique
$h\colon X\to L$ with $\pi_j\circ h = f_j$ for all~$j$.
Then $\Phi(h) = (f_j)_j$.

Naturality in~$X$ is straightforward: for $\varphi\colon X'\to X$,
$\Phi(\,h\circ\varphi\,) = (\pi_j\circ h\circ\varphi)_j$,
which is the image of $\Phi(h)$ under pre-composition with~$\varphi$.
\end{proof}

\begin{corollary}\label{cor:representable-preserves-limits}
\index{representable functor!preserves limits}
Every representable functor $\Hom(X,-)$ preserves all limits
that exist in~$\mathcal{C}$.
\end{corollary}

\begin{remark}\label{rem:colim-hom}
\index{colimit!and Hom functor}
Dually, for colimits:
\[
    \Hom_{\mathcal{C}}(\colim_j F_j,\, X)
    \;\cong\;
    \lim_j\, \Hom_{\mathcal{C}}(F_j,\, X).
\]
Note the \emph{contravariance}: the colimit on the left becomes a
limit on the right.
\end{remark}


%% ============================================================
\section{Functors and limits}\label{sec:functors-limits}
%% ============================================================

\begin{definition}[Preservation of limits]\label{def:preserve-limits}
\index{functor!preserving limits}
\index{limit!preservation}
A functor $G\colon\mathcal{C}\to\mathcal{D}$ \textbf{preserves the limit}
of a diagram $F\colon\mathcal{J}\to\mathcal{C}$ if whenever
$(L,\pi_j)$ is a limit cone for~$F$, then $(GL,\,G\pi_j)$ is a limit
cone for $G\circ F$ in~$\mathcal{D}$.

We say $G$ \textbf{preserves all limits} (or is
\textbf{continuous}\index{continuous functor}) if it preserves
the limit of every diagram that admits a limit.
\end{definition}

\begin{definition}[Reflection and creation of limits]%
\label{def:reflect-create-limits}
\index{functor!reflecting limits}
\index{functor!creating limits}
\index{limit!reflection}
\index{limit!creation}
A functor $G\colon\mathcal{C}\to\mathcal{D}$ \textbf{reflects the limit}
of $F\colon\mathcal{J}\to\mathcal{C}$ if whenever $(L,\pi_j)$ is a cone
for~$F$ such that $(GL,G\pi_j)$ is a limit cone for $G\circ F$,
then $(L,\pi_j)$ is already a limit cone for~$F$.

$G$ \textbf{creates the limit} of~$F$ if for every limit cone
$(M,\psi_j)$ for $G\circ F$ in~$\mathcal{D}$, there exists a unique
cone $(L,\pi_j)$ for~$F$ with $GL = M$ and $G\pi_j = \psi_j$,
and moreover this cone is a limit of~$F$.
\end{definition}

\begin{example}\label{ex:forgetful-creates}
\index{forgetful functor!creates limits}
The forgetful functor $U\colon\Grp\to\Set$ creates all limits.
Indeed, given a diagram of groups, one first forms the limit in~$\Set$
(using the Cartesian product construction), then shows this set carries
a unique group structure making all projections group homomorphisms.
The same is true for $\Ab\to\Set$, $\Ring\to\Set$, $\Mod_R\to\Set$, etc.
\end{example}


%% ============================================================
\section{Filtered colimits}\label{sec:filtered-colimits}
%% ============================================================

\begin{definition}[Filtered category]\label{def:filtered-category}
\index{filtered category}
A small category $\mathcal{J}$ is \textbf{filtered} if:
\begin{enumerate}[label=(\roman*)]
\item $\mathcal{J}$ is non-empty;
\item for every pair of objects $j,k\in\mathcal{J}$, there exists
      an object $\ell$ and morphisms $j\to\ell$ and $k\to\ell$;
\item for every pair of parallel morphisms $\alpha,\beta\colon j\rightrightarrows k$,
      there exists a morphism $\gamma\colon k\to\ell$ such that
      $\gamma\circ\alpha = \gamma\circ\beta$.
\end{enumerate}
\end{definition}

\begin{definition}[Filtered colimit]\label{def:filtered-colimit}
\index{filtered colimit}
\index{colimit!filtered}
A \textbf{filtered colimit} (or \textbf{direct limit}) is a colimit
of a diagram $F\colon\mathcal{J}\to\mathcal{C}$ where $\mathcal{J}$
is filtered.  We also write $\varinjlim F$.
\end{definition}

\begin{example}\label{ex:filtered-colimit-directed}
\index{directed system}
\index{directed colimit}
A partially ordered set $(I,\leq)$ that is directed
(every pair has an upper bound) is a filtered category.
A functor $F\colon I\to\mathcal{C}$ is a directed system,
and its colimit is the classical direct limit.
\end{example}

\begin{theorem}\label{thm:filtered-colim-set}
\index{filtered colimit!in Set@in $\Set$}
In~$\Set$, filtered colimits commute with finite limits.
More precisely, if $\mathcal{J}$ is filtered and $\mathcal{K}$ is finite,
then for any functor $F\colon\mathcal{J}\times\mathcal{K}\to\Set$,
\[
    \varinjlim_j\, \varprojlim_k\, F(j,k)
    \;\cong\;
    \varprojlim_k\, \varinjlim_j\, F(j,k).
\]
\end{theorem}

\begin{proof}
We sketch the proof for the key cases.  The filtered colimit
$\varinjlim_j F_j$ in~$\Set$ is computed as
$\bigl(\coprod_j F_j\bigr)/{\sim}$
where $x\in F_j$ and $y\in F_k$ satisfy $x\sim y$ iff there exist
morphisms $\alpha\colon j\to\ell$ and $\beta\colon k\to\ell$ with
$F_\alpha(x) = F_\beta(y)$.

\smallskip\noindent
\emph{Commutation with finite products:}
An element of $\varinjlim_j(F_j\times G_j)$ is an equivalence class
$[(x,y)]$ with $x\in F_j$, $y\in G_j$.  The filteredness condition~(ii)
ensures that any pair of elements from
$\varinjlim_j F_j$ and $\varinjlim_j G_j$ can be represented at a
common index, establishing a bijection with
$(\varinjlim_j F_j)\times(\varinjlim_j G_j)$.

\smallskip\noindent
\emph{Commutation with equalisers:}
Given parallel natural transformations $f,g\colon F\rightrightarrows G$,
the filteredness condition~(iii) ensures that
$\varinjlim_j\eq(f_j,g_j) \cong \eq(\varinjlim_j f_j,\,\varinjlim_j g_j)$.

Since finite products and equalisers generate all finite limits
(Theorem~\ref{thm:prod-eq-limits}), the general statement follows.
\end{proof}

\begin{example}\label{ex:filtered-colim-algebra}
\index{filtered colimit!of groups}
\index{filtered colimit!of rings}
A filtered colimit of groups is a group, a filtered colimit of rings
is a ring, etc.  For instance, $\Q = \varinjlim_{n\geq 1} \frac{1}{n}\Z$
as abelian groups, and algebraic closures can be constructed as
filtered colimits of finite extensions.
\end{example}


%% ============================================================
\section{Exercises}\label{sec:exercises-ch3}
%% ============================================================

\begin{exercise}\label{exer:terminal-limit}
\index{terminal object!as limit}
\index{initial object!as colimit}
Show that a terminal object is the limit of the empty diagram
(i.e.\ $\mathcal{J} = \varnothing$), and an initial object is the
colimit of the empty diagram.
\end{exercise}

\begin{exercise}\label{exer:equaliser-from-pullback}
\index{equaliser!from pullback}
Show that if $\mathcal{C}$ has binary products and pullbacks,
then $\mathcal{C}$ has equalisers.
\emph{Hint:} Given $f,g\colon A\rightrightarrows B$, consider the
pullback of $\langle f,g\rangle\colon A\to B\times B$ along
the diagonal $\Delta\colon B\to B\times B$.
\end{exercise}

\begin{exercise}\label{exer:product-from-pullback}
\index{product!from pullback}
Show that if $\mathcal{C}$ has a terminal object and pullbacks,
then $\mathcal{C}$ has all finite limits.
\end{exercise}

\begin{exercise}\label{exer:colimit-in-set}
\index{colimit!in Set@in $\Set$}
Show that for any small diagram $F\colon\mathcal{J}\to\Set$,
the colimit is
\[
    \varinjlim F \;=\;
    \Bigl(\coprod_{j\in\mathcal{J}} F_j\Bigr)\Big/{\sim}
\]
where $\sim$ is the equivalence relation generated by:
$x\in F_j$ is equivalent to $F_\alpha(x)\in F_k$ for every
$\alpha\colon j\to k$ in~$\mathcal{J}$.
\end{exercise}

\begin{exercise}\label{exer:limit-functor-category}
\index{limit!in functor category}
Let $\mathcal{C}$ be a complete category and $\mathcal{J}$ a small category.
Show that $\Fun(\mathcal{J},\mathcal{C})$ is also complete,
and that limits are computed ``pointwise'': for a diagram
$\Phi\colon\mathcal{K}\to\Fun(\mathcal{J},\mathcal{C})$,
\[
    (\lim_k \Phi_k)(j) \cong \lim_k(\Phi_k(j))
\]
for each $j\in\mathcal{J}$.
\end{exercise}

\begin{exercise}\label{exer:limit-preserves-mono}
\index{limit!preserves monomorphisms}
Let $\mathcal{C}$ be a category with equalisers.
Show that if $\alpha\colon F\Rightarrow G$ is a natural transformation
between diagrams $F,G\colon\mathcal{J}\to\mathcal{C}$ such that each
$\alpha_j$ is a monomorphism, and if $\lim F$ and $\lim G$ both exist,
then the induced morphism $\lim F\to\lim G$ is also a monomorphism.
\end{exercise}

\begin{exercise}\label{exer:inverse-limit-abelian-groups}
\index{inverse limit!of abelian groups}
\index{$p$-adic integers}
Let $p$ be a prime.  Consider the inverse system
$\cdots\to\Z/p^3\Z\to\Z/p^2\Z\to\Z/p\Z$
in~$\Ab$.  Show that the inverse limit is the ring of
$p$-adic integers~$\Z_p$.
\end{exercise}

\begin{exercise}\label{exer:filtered-colim-exact}
\index{filtered colimit!and exactness}
Show that filtered colimits in $\Mod_R$ (for $R$ a ring) are exact.
That is, if $0\to F\to G\to H\to 0$ is a short exact sequence
of directed systems, then
$0\to\varinjlim F\to\varinjlim G\to\varinjlim H\to 0$ is exact.
\end{exercise}

\index{limit|)}
\index{colimit|)}

\newpage

%%%%%%%%%%%%%%%%%%%%%%%%%%%%%%%%%%%%%%%%%%%%%%%%%%%%%%%%%%%%
%%           CHAPTER 4: ADJUNCTIONS                       %%
%%%%%%%%%%%%%%%%%%%%%%%%%%%%%%%%%%%%%%%%%%%%%%%%%%%%%%%%%%%%

\chapter{Adjunctions}\label{ch:adjunctions}
\minitoc

\vspace{1em}

Adjunctions are arguably the most important concept in category theory.
Saunders Mac\,Lane famously wrote that ``adjoint functors arise everywhere.''
\index{Mac Lane, Saunders}
An adjunction captures a universal relationship between two functors
going in opposite directions, generalising at a stroke the notions of
free construction, tensor--hom duality, and a host of other fundamental
correspondences throughout mathematics.

In this chapter we give three equivalent definitions of adjunctions,
prove their equivalence in detail, explore the fundamental examples,
and establish the key theorems---most notably, that left adjoints
preserve colimits and right adjoints preserve limits.  We conclude
with Freyd's adjoint functor theorems.

\index{adjunction|(}

%% ============================================================
\section{Definition via natural bijection}\label{sec:adj-hom-def}
%% ============================================================

\begin{definition}[Adjunction]\label{def:adjunction}
\index{adjunction!definition}
\index{left adjoint}
\index{right adjoint}
\index{adjunction!natural bijection}
Let $\mathcal{C}$ and $\mathcal{D}$ be categories and
$F\colon\mathcal{C}\to\mathcal{D}$,
$G\colon\mathcal{D}\to\mathcal{C}$ be functors.
We say that $F$ is \textbf{left adjoint} to~$G$
(equivalently, $G$ is \textbf{right adjoint} to~$F$),
written $F\dashv G$, if there is a natural bijection
\[
    \Phi_{A,B}\colon
    \Hom_{\mathcal{D}}(FA,\,B)
    \;\xrightarrow{\;\sim\;}\;
    \Hom_{\mathcal{C}}(A,\,GB)
\]
for all $A\in\mathcal{C}$ and $B\in\mathcal{D}$,
natural in both variables.  That is, $\Phi$ is a natural isomorphism
\[
    \Phi\colon \Hom_{\mathcal{D}}(F-,\,-)
    \;\xRightarrow{\;\sim\;}\;
    \Hom_{\mathcal{C}}(-,\,G-)
\]
of bifunctors $\mathcal{C}^{\op}\times\mathcal{D}\to\Set$.
\end{definition}

\begin{remark}\label{rem:adj-notation}
\index{adjunction!notation}
We depict an adjunction as:
\[
\begin{tikzcd}
\mathcal{C}
  \arrow[rr, bend left=25, "F", ""{name=F, below}]
& \bot &
\mathcal{D}
  \arrow[ll, bend left=25, "G", ""{name=G, above}]
  \arrow[from=F, to=G, phantom, "\dashv"]
\end{tikzcd}
\]
The symbol $\bot$ (or $\dashv$) indicates $F\dashv G$.
\end{remark}

Let us spell out what naturality means in this context.

\begin{proposition}[Naturality conditions]\label{prop:adj-naturality}
\index{adjunction!naturality}
The bijection $\Phi_{A,B}$ is natural in $A$ and $B$ if and only if:
\begin{enumerate}[label=(\roman*)]
\item \textbf{Naturality in $B$:}
For every $g\colon B\to B'$ and $f\colon FA\to B$,
\[
    \Phi_{A,B'}(g\circ f) = Gg \circ \Phi_{A,B}(f).
\]
\item \textbf{Naturality in $A$:}
For every $h\colon A'\to A$ and $f\colon FA\to B$,
\[
    \Phi_{A',B}(f\circ Fh) = \Phi_{A,B}(f)\circ h.
\]
\end{enumerate}
\end{proposition}

\begin{proof}
This is a direct translation of the naturality squares for the
bifunctor $\Hom_\mathcal{D}(F-,-) \Rightarrow \Hom_\mathcal{C}(-,G-)$.
For (i), the naturality square in~$B$ is:
\[
\begin{tikzcd}
\Hom(FA,B) \arrow[r, "\Phi_{A,B}"] \arrow[d, "g\circ -"'] &
\Hom(A,GB) \arrow[d, "Gg\circ -"] \\
\Hom(FA,B') \arrow[r, "\Phi_{A,B'}"'] & \Hom(A,GB')
\end{tikzcd}
\]
Commutativity gives $Gg\circ\Phi_{A,B}(f) = \Phi_{A,B'}(g\circ f)$.
The argument for (ii) is analogous.
\end{proof}


%% ============================================================
\section{Unit and counit}\label{sec:unit-counit}
%% ============================================================

\begin{definition}[Unit and counit]\label{def:unit-counit}
\index{unit of adjunction}
\index{counit of adjunction}
\index{adjunction!unit}
\index{adjunction!counit}
Given an adjunction $F\dashv G$ with bijection~$\Phi$, define:
\begin{itemize}
\item The \textbf{unit} $\eta\colon \Id_{\mathcal{C}}\Rightarrow GF$
    by $\eta_A = \Phi_{A,FA}(\id_{FA})$ for each $A\in\mathcal{C}$.
\item The \textbf{counit} $\varepsilon\colon FG\Rightarrow\Id_{\mathcal{D}}$
    by $\varepsilon_B = \Phi^{-1}_{GB,B}(\id_{GB})$ for each
    $B\in\mathcal{D}$.
\end{itemize}
\end{definition}

\begin{proposition}\label{prop:unit-counit-nat}
\index{unit of adjunction!naturality}
\index{counit of adjunction!naturality}
The unit $\eta$ and counit $\varepsilon$ are natural transformations.
Moreover, the bijection $\Phi$ can be recovered from them:
\[
    \Phi_{A,B}(f\colon FA\to B) = Gf\circ\eta_A,
    \qquad
    \Phi^{-1}_{A,B}(g\colon A\to GB) = \varepsilon_B\circ Fg.
\]
\end{proposition}

\begin{proof}
\textbf{Recovery formula.}
For $f\colon FA\to B$, naturality of~$\Phi$ in~$B$ gives
\[
    \Phi_{A,B}(f) = \Phi_{A,B}(f\circ\id_{FA})
    = Gf\circ\Phi_{A,FA}(\id_{FA}) = Gf\circ\eta_A.
\]
Similarly, for $g\colon A\to GB$, naturality in~$A$ gives
\[
    \Phi^{-1}_{A,B}(g)
    = \Phi^{-1}_{A,B}(\id_{GB}\circ g)
    = \Phi^{-1}_{GB,B}(\id_{GB})\circ Fg
    = \varepsilon_B\circ Fg.
\]

\textbf{Naturality of $\eta$.}
We must show that for $h\colon A\to A'$, $\eta_{A'}\circ h = GFh\circ\eta_A$:
\[
\begin{tikzcd}
A \arrow[r, "h"] \arrow[d, "\eta_A"'] & A' \arrow[d, "\eta_{A'}"] \\
GFA \arrow[r, "GFh"'] & GFA'
\end{tikzcd}
\]
We have $\eta_{A'}\circ h = \Phi_{A',FA'}(\id_{FA'})\circ h
= \Phi_{A,FA'}(\id_{FA'}\circ Fh) = \Phi_{A,FA'}(Fh)
= GFh\circ\eta_A$, using naturality in~$A$ for the second
equality and the recovery formula for the last.

\textbf{Naturality of $\varepsilon$} is proved dually.
\end{proof}


%% ============================================================
\section{Triangle identities}\label{sec:triangle-identities}
%% ============================================================

\begin{theorem}[Triangle identities]\label{thm:triangle-identities}
\index{triangle identities}
\index{zig-zag identities}
\index{adjunction!triangle identities}
Let $F\dashv G$ with unit $\eta$ and counit~$\varepsilon$.
Then the following two composites are identities:
\begin{enumerate}[label=(\alph*)]
\item $\varepsilon F \circ F\eta = \id_F$, i.e.\ for every $A$:
\[
\begin{tikzcd}
FA \arrow[r, "F\eta_A"] \arrow[dr, equal] & FGFA \arrow[d, "\varepsilon_{FA}"] \\
& FA
\end{tikzcd}
\]
\item $G\varepsilon \circ \eta G = \id_G$, i.e.\ for every $B$:
\[
\begin{tikzcd}
GB \arrow[r, "\eta_{GB}"] \arrow[dr, equal] & GFGB \arrow[d, "G\varepsilon_B"] \\
& GB
\end{tikzcd}
\]
\end{enumerate}
These are called the \textbf{triangle identities} (or \textbf{zig-zag
identities}) because of their shape in the following diagram:
\[
\begin{tikzcd}[column sep=large]
F \arrow[r, Rightarrow, "F\eta"] \arrow[dr, equal]
  & FGF \arrow[d, Rightarrow, "\varepsilon F"] &
G \arrow[r, Rightarrow, "\eta G"] \arrow[dr, equal]
  & GFG \arrow[d, Rightarrow, "G\varepsilon"] \\
& F & & G
\end{tikzcd}
\]
\end{theorem}

\begin{proof}
\textbf{(a):}
We must show $\varepsilon_{FA}\circ F\eta_A = \id_{FA}$.
By the recovery formula (Proposition~\ref{prop:unit-counit-nat}),
$\Phi^{-1}_{A,FA}(\eta_A) = \varepsilon_{FA}\circ F\eta_A$.
But $\eta_A = \Phi_{A,FA}(\id_{FA})$, so
$\Phi^{-1}_{A,FA}(\eta_A) = \id_{FA}$.
Hence $\varepsilon_{FA}\circ F\eta_A = \id_{FA}$.

\medskip
\textbf{(b):}
We must show $G\varepsilon_B\circ\eta_{GB} = \id_{GB}$.
By the recovery formula,
$\Phi_{GB,B}(\varepsilon_B) = G\varepsilon_B\circ\eta_{GB}$.
But $\varepsilon_B = \Phi^{-1}_{GB,B}(\id_{GB})$, so
$\Phi_{GB,B}(\varepsilon_B) = \id_{GB}$.
Hence $G\varepsilon_B\circ\eta_{GB} = \id_{GB}$.
\end{proof}


%% ============================================================
\section{Equivalence of definitions}\label{sec:adj-equivalence}
%% ============================================================

We now state and prove the equivalence of three different ways
of specifying an adjunction.

\begin{theorem}[Three equivalent definitions of adjunction]%
\label{thm:adj-equiv-defs}
\index{adjunction!equivalent definitions}
Let $F\colon\mathcal{C}\to\mathcal{D}$ and
$G\colon\mathcal{D}\to\mathcal{C}$ be functors.
The following data are in bijective correspondence:
\begin{enumerate}[label=\textbf{(\Alph*)}]
\item\label{item:adj-hom}
    A natural isomorphism
    $\Phi\colon\Hom_\mathcal{D}(F-,-)\xRightarrow{\sim}\Hom_\mathcal{C}(-,G-)$.

\item\label{item:adj-unit-counit}
    Natural transformations $\eta\colon\Id_\mathcal{C}\Rightarrow GF$
    and $\varepsilon\colon FG\Rightarrow\Id_\mathcal{D}$ satisfying
    the triangle identities:
    \[
        \varepsilon F\circ F\eta = \id_F,
        \qquad
        G\varepsilon\circ\eta G = \id_G.
    \]

\item\label{item:adj-universal}
    For each $A\in\mathcal{C}$, an object $FA\in\mathcal{D}$ and
    a morphism $\eta_A\colon A\to GFA$ that is \textbf{universal}:
    for every $B\in\mathcal{D}$ and $g\colon A\to GB$, there exists
    a unique $f\colon FA\to B$ with $Gf\circ\eta_A = g$:
    \[
    \begin{tikzcd}
    A \arrow[r, "\eta_A"] \arrow[dr, "g"'] & GFA \arrow[d, "Gf"] \\
    & GB
    \end{tikzcd}
    \]
\end{enumerate}
\end{theorem}

\begin{proof}
\textbf{\ref{item:adj-hom} $\Rightarrow$ \ref{item:adj-unit-counit}:}
Given $\Phi$, define $\eta_A = \Phi_{A,FA}(\id_{FA})$ and
$\varepsilon_B = \Phi^{-1}_{GB,B}(\id_{GB})$.
By Proposition~\ref{prop:unit-counit-nat}, these are natural transformations.
The triangle identities are proved in Theorem~\ref{thm:triangle-identities}.

\medskip
\textbf{\ref{item:adj-unit-counit} $\Rightarrow$ \ref{item:adj-hom}:}
Given $(\eta,\varepsilon)$ satisfying the triangle identities, define
\[
    \Phi_{A,B}(f) = Gf\circ\eta_A, \qquad
    \Psi_{A,B}(g) = \varepsilon_B\circ Fg.
\]
We verify $\Psi\circ\Phi = \id$ and $\Phi\circ\Psi = \id$.

For $f\colon FA\to B$:
\begin{align*}
    \Psi(\Phi(f)) &= \Psi(Gf\circ\eta_A)
    = \varepsilon_B\circ F(Gf\circ\eta_A) \\
    &= \varepsilon_B\circ FGf\circ F\eta_A
    = f\circ\varepsilon_{FA}\circ F\eta_A \quad\text{(naturality of $\varepsilon$)}\\
    &= f\circ\id_{FA} = f \quad\text{(triangle identity (a))}.
\end{align*}

For $g\colon A\to GB$:
\begin{align*}
    \Phi(\Psi(g)) &= \Phi(\varepsilon_B\circ Fg)
    = G(\varepsilon_B\circ Fg)\circ\eta_A \\
    &= G\varepsilon_B\circ GFg\circ\eta_A
    = G\varepsilon_B\circ\eta_{GB}\circ g \quad\text{(naturality of $\eta$)}\\
    &= \id_{GB}\circ g = g \quad\text{(triangle identity (b))}.
\end{align*}

Naturality of $\Phi$ in both variables follows from the naturality of
$\eta$ and the functoriality of~$G$.

\medskip
\textbf{\ref{item:adj-hom} $\Leftrightarrow$ \ref{item:adj-universal}:}
Given~$\Phi$, setting $\eta_A = \Phi_{A,FA}(\id_{FA})$ provides the
universal morphism.  The universality statement is precisely the
bijectivity of~$\Phi_{A,B}$: given $g\colon A\to GB$, the unique
$f$ with $Gf\circ\eta_A = g$ is $f = \Phi^{-1}_{A,B}(g)$.

Conversely, given universal arrows $(\eta_A)$, define
$\Phi_{A,B}(f) = Gf\circ\eta_A$.  The universality ensures this is
a bijection.  Naturality follows from the naturality of the collection
$(\eta_A)_A$ (which itself follows from universality applied to the
composites $GFh\circ\eta_A$).
\end{proof}


%% ============================================================
\section{Fundamental examples}\label{sec:adj-examples}
%% ============================================================

\subsection{Free--forgetful adjunctions}

\begin{example}[Free group]\label{ex:free-group}
\index{adjunction!free-forgetful}
\index{free group}
\index{free functor}
\index{forgetful functor}
Let $U\colon\Grp\to\Set$ be the forgetful functor and
$F\colon\Set\to\Grp$ the free group functor.
Then $F\dashv U$:
\[
\begin{tikzcd}
\Set
  \arrow[rr, bend left=25, "F", ""{name=F, below}]
& \bot &
\Grp
  \arrow[ll, bend left=25, "U", ""{name=G, above}]
\end{tikzcd}
\]
The adjunction bijection is
\[
    \Hom_\Grp(F(S),\,G) \;\cong\; \Hom_\Set(S,\,U(G)):
\]
a group homomorphism from the free group on~$S$ is determined by
where the generators go, i.e.\ by a function $S\to U(G)$.

The unit $\eta_S\colon S\to UF(S)$ sends each element to the
corresponding generator in the free group.
The counit $\varepsilon_G\colon FU(G)\to G$ sends a word in the
``generators'' $U(G)$ to its evaluation in~$G$.
\end{example}

\begin{example}[Free abelian group]\label{ex:free-abelian}
\index{free abelian group}
The free abelian group functor $\Z[-]\colon\Set\to\Ab$ is left adjoint
to the forgetful functor $U\colon\Ab\to\Set$:
\[
    \Hom_\Ab(\Z[S],\,A) \;\cong\; \Hom_\Set(S,\,U(A)).
\]
Here $\Z[S] = \bigoplus_{s\in S} \Z$ is the free abelian group on~$S$.
\end{example}

\begin{example}[Free $R$-module]\label{ex:free-module}
\index{free module}
For a ring~$R$, the free module functor $R[-]\colon\Set\to\Mod_R$
is left adjoint to the forgetful functor.
\end{example}

\begin{example}[Free--forgetful between $\Ab$ and $\Grp$]%
\label{ex:abelianisation}
\index{abelianisation}
The inclusion functor $\iota\colon\Ab\hookrightarrow\Grp$ has a
left adjoint: the abelianisation $\mathrm{ab}\colon\Grp\to\Ab$
sending $G$ to $G/[G,G]$.
\[
    \Hom_\Ab(G/[G,G],\,A) \;\cong\; \Hom_\Grp(G,\,\iota A).
\]
\end{example}

\subsection{Tensor--Hom adjunction}

\begin{example}[Tensor--Hom]\label{ex:tensor-hom}
\index{adjunction!tensor-Hom}
\index{tensor product!as left adjoint}
\index{Hom!as right adjoint}
Let $R$ be a commutative ring and $M$ an $R$-module.
Then $(-\otimes_R M) \dashv \Hom_R(M,-)$:
\[
    \Hom_R(N\otimes_R M,\, P) \;\cong\; \Hom_R(N,\,\Hom_R(M,P))
\]
naturally in $N$ and~$P$.  This is the universal property of the
tensor product.
\end{example}

\subsection{Product--Internal Hom (Cartesian closed categories)}

\begin{example}[Cartesian closed]\label{ex:cartesian-closed}
\index{adjunction!product-internal Hom}
\index{Cartesian closed category}
\index{exponential object}
In a Cartesian closed category (e.g.\ $\Set$), for each object~$B$
the product functor $(-\times B)$ has a right adjoint, the
\textbf{internal Hom} (or \textbf{exponential}) $B^{(-)}$ or $[B,-]$:
\[
    \Hom(A\times B,\,C) \;\cong\; \Hom(A,\,C^B).
\]
In $\Set$, $C^B = \{f\colon B\to C\}$ is the set of functions.
This is the \textbf{currying} correspondence in computer science.
\index{currying}
\end{example}

\subsection{Direct and inverse image (sheaves)}

\begin{example}[Direct and inverse image]\label{ex:sheaf-adjunction}
\index{adjunction!direct and inverse image}
\index{sheaves!adjunction}
\index{inverse image functor}
\index{direct image functor}
Let $f\colon X\to Y$ be a continuous map of topological spaces.
There is an adjunction
\[
    f^* \dashv f_*
\]
between the inverse image functor $f^*\colon\mathsf{Sh}(Y)\to\mathsf{Sh}(X)$
and the direct image functor $f_*\colon\mathsf{Sh}(X)\to\mathsf{Sh}(Y)$
on categories of sheaves.  This adjunction is fundamental in algebraic
geometry and sheaf theory.
\end{example}

\subsection{Suspension--loops adjunction}

\begin{example}[Suspension and loops]\label{ex:suspension-loops}
\index{adjunction!suspension-loops}
\index{suspension functor}
\index{loop space functor}
In the category $\Top_*$ of pointed topological spaces, the
(reduced) suspension functor $\Sigma$ is left adjoint to the
loops functor~$\Omega$:
\[
    \Sigma \dashv \Omega, \qquad
    [\Sigma X,\, Y]_* \;\cong\; [X,\,\Omega Y]_*
\]
where $[-,-]_*$ denotes pointed homotopy classes.
This adjunction underlies the theory of stable homotopy.
\end{example}

\begin{example}[Galois connections]\label{ex:galois-connection}
\index{Galois connection}
\index{adjunction!as Galois connection}
If $\mathcal{C}$ and $\mathcal{D}$ are posets (viewed as categories),
then an adjunction $F\dashv G$ is the same as a \textbf{Galois connection}:
$F(a)\leq b$ if and only if $a\leq G(b)$.
\end{example}


%% ============================================================
\section{Left adjoints preserve colimits}\label{sec:ladj-colim}
%% ============================================================

\begin{theorem}[Left adjoints preserve colimits]%
\label{thm:ladj-preserve-colim}
\index{left adjoint!preserves colimits}
\index{colimit!preserved by left adjoint}
\index{adjunction!and colimits}
If $F\dashv G$ with $F\colon\mathcal{C}\to\mathcal{D}$,
then $F$ preserves all colimits that exist in~$\mathcal{C}$.
\end{theorem}

\begin{proof}
Let $D\colon\mathcal{J}\to\mathcal{C}$ be a diagram with colimit
$(\varinjlim D,\,\iota_j)$.  We must show that
$(F(\varinjlim D),\,F\iota_j)$ is a colimit of $F\circ D$.

Let $(N,\psi_j)$ be a cocone under $F\circ D$ in~$\mathcal{D}$,
so $\psi_j\colon FD_j\to N$ with compatibility.
For each~$j$, applying the adjunction bijection gives
\[
    \tilde\psi_j = \Phi_{D_j,N}(\psi_j) \colon D_j \to GN.
\]

\textbf{Claim:} $(GN,\tilde\psi_j)$ is a cocone under~$D$ in~$\mathcal{C}$.
Indeed, for $\alpha\colon j\to k$ in~$\mathcal{J}$:
\begin{align*}
\tilde\psi_k\circ D_\alpha
&= \Phi_{D_k,N}(\psi_k)\circ D_\alpha \\
&= \Phi_{D_j,N}(\psi_k\circ FD_\alpha)
    \quad\text{(naturality of $\Phi$ in $A$)} \\
&= \Phi_{D_j,N}(\psi_j) = \tilde\psi_j
    \quad\text{(cocone condition for $\psi$)}.
\end{align*}

By the universal property of $\varinjlim D$, there is a unique
$\tilde u\colon\varinjlim D\to GN$ with
$\tilde u\circ\iota_j = \tilde\psi_j$ for all~$j$.

Now set $u = \Phi^{-1}_{\varinjlim D,\,N}(\tilde u)\colon F(\varinjlim D)\to N$.
Then for each~$j$, using naturality of~$\Phi^{-1}$ in~$A$:
\[
    u\circ F\iota_j
    = \Phi^{-1}_{D_j,N}(\tilde u\circ\iota_j)
    = \Phi^{-1}_{D_j,N}(\tilde\psi_j)
    = \psi_j.
\]
This gives existence.

\textbf{Uniqueness:}
If $u'\colon F(\varinjlim D)\to N$ also satisfies
$u'\circ F\iota_j = \psi_j$, then
$\Phi_{\varinjlim D,N}(u')\circ\iota_j = \tilde\psi_j$ for all~$j$,
so by uniqueness for $\varinjlim D$,
$\Phi(u') = \tilde u = \Phi(u)$, hence $u' = u$.
\[
\begin{tikzcd}
FD_j \arrow[r, "F\iota_j"] \arrow[dr, "\psi_j"']
& F(\varinjlim D) \arrow[d, dashed, "\exists!\, u"] \\
& N
\end{tikzcd}
\]
\end{proof}


%% ============================================================
\section{Right adjoints preserve limits}\label{sec:radj-lim}
%% ============================================================

\begin{theorem}[Right adjoints preserve limits]%
\label{thm:radj-preserve-lim}
\index{right adjoint!preserves limits}
\index{limit!preserved by right adjoint}
\index{adjunction!and limits}
If $F\dashv G$ with $G\colon\mathcal{D}\to\mathcal{C}$,
then $G$ preserves all limits that exist in~$\mathcal{D}$.
\end{theorem}

\begin{proof}
Let $D\colon\mathcal{J}\to\mathcal{D}$ be a diagram with limit
$(L,\pi_j)$, where $\pi_j\colon L\to D_j$.
We show $(GL,G\pi_j)$ is a limit of $G\circ D$.

Let $(N,\psi_j)$ be a cone over $G\circ D$ in~$\mathcal{C}$,
so $\psi_j\colon N\to GD_j$.
By the adjunction, for each~$j$ let
$\tilde\psi_j = \Phi^{-1}_{N,D_j}(\psi_j)\colon FN\to D_j$.

\textbf{Claim:} $(FN,\tilde\psi_j)$ is a cone over~$D$.
For $\alpha\colon j\to k$:
\begin{align*}
D_\alpha\circ\tilde\psi_j
&= D_\alpha\circ\Phi^{-1}_{N,D_j}(\psi_j) \\
&= \Phi^{-1}_{N,D_k}(GD_\alpha\circ\psi_j)
  \quad\text{(naturality of $\Phi^{-1}$ in $B$)} \\
&= \Phi^{-1}_{N,D_k}(\psi_k) = \tilde\psi_k
  \quad\text{(cone condition for $\psi$)}.
\end{align*}

By the universal property of~$L$, there is a unique
$\tilde v\colon FN\to L$ with $\pi_j\circ\tilde v = \tilde\psi_j$.
Set $v = \Phi_{N,L}(\tilde v)\colon N\to GL$.  Then
\[
    G\pi_j\circ v = G\pi_j\circ\Phi_{N,L}(\tilde v)
    = \Phi_{N,D_j}(\pi_j\circ\tilde v)
    = \Phi_{N,D_j}(\tilde\psi_j) = \psi_j,
\]
using naturality of~$\Phi$ in~$B$.

\textbf{Uniqueness} follows as before from the bijectivity of~$\Phi$
and uniqueness for the limit~$L$.
\[
\begin{tikzcd}
N \arrow[dr, "\psi_j"'] \arrow[rr, dashed, "\exists!\, v"]
& & GL \arrow[dl, "G\pi_j"] \\
& GD_j &
\end{tikzcd}
\]
\end{proof}

\begin{corollary}\label{cor:limits-detect-adjoints}
\index{adjunction!and preservation of (co)limits}
If a functor $G\colon\mathcal{D}\to\mathcal{C}$ has a left adjoint,
then $G$ preserves all limits.
If $F\colon\mathcal{C}\to\mathcal{D}$ has a right adjoint,
then $F$ preserves all colimits.
\end{corollary}

\begin{remark}\label{rem:converse-not-true}
The converse does not hold in general: preserving limits is necessary
but not sufficient for having a left adjoint.  The adjoint functor
theorems (below) provide sufficient conditions.
\end{remark}

\begin{example}\label{ex:forgetful-preserves-limits}
\index{forgetful functor!preserves limits}
Since the forgetful functor $U\colon\Grp\to\Set$ has a left adjoint
(the free group functor), it preserves all limits.  This explains why
limits of groups are computed ``on underlying sets'':
the underlying set of a product of groups is the product of their
underlying sets, etc.
\end{example}


%% ============================================================
\section{Uniqueness of adjoints}\label{sec:adj-unique}
%% ============================================================

\begin{theorem}[Uniqueness of adjoints]\label{thm:adj-unique}
\index{adjunction!uniqueness}
\index{right adjoint!uniqueness}
\index{left adjoint!uniqueness}
If $F\dashv G$ and $F\dashv G'$, then $G\cong G'$ via a natural
isomorphism.  Similarly, if $F\dashv G$ and $F'\dashv G$, then
$F\cong F'$.
\end{theorem}

\begin{proof}
Suppose $F\dashv G$ and $F\dashv G'$.  For every $B\in\mathcal{D}$,
\[
    \Hom_\mathcal{C}(-,GB)
    \cong \Hom_\mathcal{D}(F-,B)
    \cong \Hom_\mathcal{C}(-,G'B)
\]
naturally.  By the Yoneda lemma, a natural isomorphism
$\Hom(-,GB)\cong\Hom(-,G'B)$ implies $GB\cong G'B$,
and the naturality in~$B$ gives a natural isomorphism $G\cong G'$.

The uniqueness of left adjoints follows by a dual argument
(using the covariant Yoneda embedding).
\end{proof}


%% ============================================================
\section{Freyd's Adjoint Functor Theorems}\label{sec:GAFT}
%% ============================================================

The adjoint functor theorems give powerful sufficient conditions
for a functor to have an adjoint.  They are among the deepest
results of basic category theory.

\begin{definition}[Solution set condition]\label{def:solution-set}
\index{solution set condition}
A functor $G\colon\mathcal{D}\to\mathcal{C}$ satisfies the
\textbf{solution set condition} if for each object $A\in\mathcal{C}$,
there exists a \emph{set} (not a proper class) of morphisms
$\{f_i\colon A\to GB_i\}_{i\in I}$ such that every morphism
$g\colon A\to GB$ factors as $g = Gh\circ f_i$ for some $i\in I$
and some $h\colon B_i\to B$:
\[
\begin{tikzcd}
A \arrow[r, "f_i"] \arrow[dr, "g"'] & GB_i \arrow[d, "Gh"] \\
& GB
\end{tikzcd}
\]
\end{definition}

\begin{theorem}[General Adjoint Functor Theorem (Freyd)]%
\label{thm:GAFT}
\index{General Adjoint Functor Theorem}
\index{GAFT}
\index{Freyd's Adjoint Functor Theorem}
\index{adjoint functor theorem!general}
Let $\mathcal{D}$ be a complete, locally small category and let
$G\colon\mathcal{D}\to\mathcal{C}$ be a functor.
Then $G$ has a left adjoint if and only if:
\begin{enumerate}[label=(\roman*)]
\item $G$ preserves all (small) limits; and
\item $G$ satisfies the solution set condition.
\end{enumerate}
\end{theorem}

\begin{proof}
\textbf{Necessity:}
If $G$ has a left adjoint~$F$, then $G$ preserves limits by
Theorem~\ref{thm:radj-preserve-lim}.
The solution set condition is satisfied with the singleton
$\{\eta_A\colon A\to GFA\}$ for each~$A$ (where $\eta$ is the unit):
any $g\colon A\to GB$ equals $Gf\circ\eta_A$ where
$f = \Phi^{-1}(g)\colon FA\to B$.

\medskip
\textbf{Sufficiency:}
We construct, for each $A\in\mathcal{C}$, a universal arrow
$\eta_A\colon A\to GFA$ (Definition~\ref{def:adjunction},
form~\ref{item:adj-universal}).

\medskip\noindent
\emph{Step 1: Initial construction.}
Let $\{f_i\colon A\to GB_i\}_{i\in I}$ be a solution set for~$A$.
Form the product $P = \prod_{i\in I} B_i$ in~$\mathcal{D}$
(which exists since $\mathcal{D}$ is complete).
Since $G$ preserves limits, $GP \cong \prod_i GB_i$,
so the $f_i$ assemble into a single morphism
$f\colon A\to GP$.

\medskip\noindent
\emph{Step 2: Equalising.}
Consider the set of all endomorphisms $e$ of~$P$ in~$\mathcal{D}$
such that $Ge\circ f = f$.  Let $E$ be the joint equaliser of all
such endomorphisms (an intersection of equalisers, which exists by
completeness).  Let $m\colon E\to P$ be the canonical morphism
and let $\eta_A = Gm^{-1}\circ f'\colon A\to GE$
be the factored morphism.  More precisely, $f$ factors through $GE$
via a unique $\eta_A\colon A\to GE$ with $Gm\circ\eta_A = f$.

\medskip\noindent
\emph{Step 3: Verification of universality.}
Given $g\colon A\to GB$, the solution set condition gives
$g = Gh\circ f_i$ for some~$i$ and some $h\colon B_i\to B$.
Define $h'\colon P\to B$ extending~$h$ via the product projection.
Then $Gh'\circ f = Gh'\circ Gm\circ\eta_A$,
and one verifies that $h'$ factors through~$E$
(by the equaliser condition and the fact that $G$ preserves the
relevant limits), yielding the desired factorisation
$g = G\bar h\circ\eta_A$ for a unique $\bar h\colon E\to B$.

\medskip\noindent
\emph{Uniqueness} of $\bar h$ requires a further equaliser argument:
if $\bar h,\bar h'\colon E\to B$ both satisfy $G\bar h\circ\eta_A
= G\bar h'\circ\eta_A$, one shows that the equaliser of $\bar h$
and~$\bar h'$ is all of~$E$, using the fact that $G$ preserves
this equaliser.

Setting $FA = E$ for each $A$ and using the universal arrows to define
$F$ on morphisms, one obtains the left adjoint
$F\colon\mathcal{C}\to\mathcal{D}$.
\end{proof}

\begin{theorem}[Special Adjoint Functor Theorem]%
\label{thm:SAFT}
\index{Special Adjoint Functor Theorem}
\index{SAFT}
\index{adjoint functor theorem!special}
\index{cogenerating set}
\index{well-powered category}
Let $\mathcal{D}$ be a complete, locally small, well-powered
category with a cogenerating set.
Then a functor $G\colon\mathcal{D}\to\mathcal{C}$ has a left
adjoint if and only if $G$ preserves all (small) limits.
\end{theorem}

In other words, the solution set condition is automatic when $\mathcal{D}$
has a cogenerating set and is well-powered.

\begin{remark}\label{rem:SAFT-applications}
\index{Special Adjoint Functor Theorem!applications}
The SAFT applies notably to:
\begin{itemize}
\item $\mathcal{D} = \Set$ (cogenerating set: $\{1, 2\}$ where
      $2 = \{0,1\}$);
\item $\mathcal{D} = \Ab$ (cogenerating set: $\{\Q/\Z\}$);
\item $\mathcal{D} = \Mod_R$ for $R$ Noetherian
      (any injective cogenerator suffices).
\end{itemize}
For instance, the SAFT immediately implies that any limit-preserving
functor $G\colon\Ab\to\mathcal{C}$ from the abelian category~$\Ab$
has a left adjoint, since $\Ab$ is complete, locally small,
well-powered, and has $\Q/\Z$ as a cogenerator.
\end{remark}

\begin{example}[Application of GAFT]\label{ex:GAFT-application}
\index{General Adjoint Functor Theorem!application}
\index{Stone--\v{C}ech compactification}
Consider the inclusion functor $G\colon\mathsf{CompHaus}\hookrightarrow\Top$
from compact Hausdorff spaces to all topological spaces.
Since $G$ preserves all limits (closed subspaces of compact Hausdorff
spaces are compact Hausdorff) and the solution set condition is verified
by cardinality arguments, the GAFT gives a left adjoint
$\beta\colon\Top\to\mathsf{CompHaus}$, the
\textbf{Stone--\v{C}ech compactification}.
\end{example}


%% ============================================================
\section{Composition of adjunctions}\label{sec:adj-composition}
%% ============================================================

\begin{proposition}[Composition of adjunctions]\label{prop:adj-compose}
\index{adjunction!composition}
If $F_1\dashv G_1$ between $\mathcal{C}$ and $\mathcal{D}$,
and $F_2\dashv G_2$ between $\mathcal{D}$ and $\mathcal{E}$, then
$F_2 F_1 \dashv G_1 G_2$:
\[
\begin{tikzcd}
\mathcal{C}
  \arrow[r, bend left, "F_1"]
& \mathcal{D}
  \arrow[l, bend left, "G_1"]
  \arrow[r, bend left, "F_2"]
& \mathcal{E}
  \arrow[l, bend left, "G_2"]
\end{tikzcd}
\qquad\Longrightarrow\qquad
\begin{tikzcd}
\mathcal{C}
  \arrow[rr, bend left, "F_2 F_1"]
& &
\mathcal{E}
  \arrow[ll, bend left, "G_1 G_2"]
\end{tikzcd}
\]
\end{proposition}

\begin{proof}
For $A\in\mathcal{C}$ and $E\in\mathcal{E}$:
\[
    \Hom_\mathcal{E}(F_2 F_1 A,\, E)
    \cong \Hom_\mathcal{D}(F_1 A,\, G_2 E)
    \cong \Hom_\mathcal{C}(A,\, G_1 G_2 E),
\]
naturally in $A$ and~$E$.  The composite of natural isomorphisms is
a natural isomorphism.
\end{proof}


%% ============================================================
\section{Adjunctions and equivalences}\label{sec:adj-equiv}
%% ============================================================

\begin{proposition}\label{prop:equiv-is-adjunction}
\index{equivalence of categories!and adjunctions}
\index{adjunction!and equivalence}
Every equivalence of categories gives rise to an adjunction.
Specifically, if $F\colon\mathcal{C}\to\mathcal{D}$ is part of an
equivalence $(F,G,\eta,\varepsilon)$ with $\eta\colon\Id\xRightarrow{\sim}GF$
and $\varepsilon\colon FG\xRightarrow{\sim}\Id$ natural isomorphisms,
then by modifying $\varepsilon$ if necessary, one obtains $F\dashv G$.
\end{proposition}

\begin{proof}
Given the equivalence, set $\varepsilon' = \varepsilon \circ F(\eta^{-1})_{G}
\circ (F\eta G)^{-1}$---more precisely, define
$\varepsilon'_B = \varepsilon_B \circ F(G\varepsilon_B)^{-1} \circ F\eta_{GB}$?
The standard trick is: set
\[
    \varepsilon'_B = \varepsilon_B \circ F(\eta^{-1}_{GB}) \circ F\eta_{GB}
\]
which simplifies, but let us use a cleaner approach.
Define the new counit as
$\varepsilon'_B = \varepsilon_B$ and check whether the triangle
identity holds up to adjusting signs.

In fact, the cleaner argument is: define $\Phi_{A,B}(f) = Gf\circ\eta_A$.
Since $\eta_A$ is an isomorphism, $\Phi_{A,B}$ is a bijection
(given $g\colon A\to GB$, set $f = \varepsilon_B\circ F(\eta_A^{-1}\circ
\cdots)$---the details require care).

Alternatively, observe that $(F,G,\eta,\varepsilon')$ with
$\varepsilon'_B = \varepsilon_B\circ F(G\varepsilon_B^{-1}\circ\eta_{GB})^{-1}
\circ F\eta_{GB}$ satisfies the triangle identities
after suitable simplification using the fact that $\eta$ and $\varepsilon$
are natural isomorphisms.  We omit the (routine but notationally heavy)
verification.
\end{proof}

\begin{remark}\label{rem:adjoint-equivalence}
\index{adjoint equivalence}
An \textbf{adjoint equivalence} is an adjunction $F\dashv G$ in which
both $\eta$ and $\varepsilon$ are natural isomorphisms.
Every equivalence of categories can be promoted to an adjoint equivalence.
\end{remark}


%% ============================================================
\section{Adjunctions via initial objects}\label{sec:adj-comma}
%% ============================================================

\begin{remark}\label{rem:adj-comma-cat}
\index{comma category}
\index{adjunction!via comma category}
Yet another characterisation: $G\colon\mathcal{D}\to\mathcal{C}$ has a
left adjoint if and only if for every $A\in\mathcal{C}$, the
\textbf{comma category} $(A\downarrow G)$ has an initial object.
The initial object is the pair $(FA,\eta_A\colon A\to GFA)$.
Recall that objects of $(A\downarrow G)$ are pairs $(B,\,f\colon A\to GB)$,
and a morphism $(B,f)\to(B',f')$ is an $h\colon B\to B'$ with
$Gh\circ f = f'$.
\end{remark}


%% ============================================================
\section{Exercises}\label{sec:exercises-ch4}
%% ============================================================

\begin{exercise}\label{exer:adj-check-triangle}
\index{triangle identities!verification}
Verify the triangle identities directly for the free--forgetful adjunction
$F\dashv U$ between $\Set$ and $\Grp$ (Example~\ref{ex:free-group}).
\end{exercise}

\begin{exercise}\label{exer:adj-poset}
\index{Galois connection!examples}
Let $f\colon X\to Y$ be a function between sets.  Define
$f_*\colon\mathcal{P}(X)\to\mathcal{P}(Y)$ by $f_*(S)=f(S)$
(direct image) and $f^*\colon\mathcal{P}(Y)\to\mathcal{P}(X)$
by $f^*(T)=f^{-1}(T)$ (inverse image), where $\mathcal{P}$ denotes
power set ordered by inclusion.
Show that $f_*\dashv f^*$ as a Galois connection.
Also show that $f^*$ has a right adjoint $f_!\colon\mathcal{P}(X)\to\mathcal{P}(Y)$
given by $f_!(S) = \{y\in Y : f^{-1}(\{y\})\subseteq S\}$.
\end{exercise}

\begin{exercise}\label{exer:adj-tensor-hom-triangle}
\index{adjunction!tensor-Hom!unit and counit}
For the tensor--Hom adjunction $(-\otimes_R M)\dashv\Hom_R(M,-)$
(Example~\ref{ex:tensor-hom}):
\begin{enumerate}[label=(\alph*)]
\item Describe the unit $\eta_N\colon N\to\Hom_R(M,\,N\otimes_R M)$
    explicitly.
\item Describe the counit
    $\varepsilon_P\colon\Hom_R(M,P)\otimes_R M\to P$ explicitly.
\item Verify the triangle identities.
\end{enumerate}
\end{exercise}

\begin{exercise}\label{exer:ladj-preserves-colim-examples}
\index{left adjoint!preserves colimits!examples}
Use the fact that left adjoints preserve colimits to prove:
\begin{enumerate}[label=(\alph*)]
\item The free group functor $F\colon\Set\to\Grp$ preserves coproducts.
    Conclude that $F(S_1\amalg S_2)\cong F(S_1)*F(S_2)$ (free product).
\item The tensor product $-\otimes_R M$ preserves direct sums:
    $(\bigoplus_i N_i)\otimes_R M \cong \bigoplus_i(N_i\otimes_R M)$.
\end{enumerate}
\end{exercise}

\begin{exercise}\label{exer:radj-preserves-lim-examples}
\index{right adjoint!preserves limits!examples}
Use the fact that right adjoints preserve limits to prove:
\begin{enumerate}[label=(\alph*)]
\item The forgetful functor $U\colon\Grp\to\Set$ preserves products.
\item $\Hom_R(M,\prod_i P_i)\cong\prod_i\Hom_R(M,P_i)$.
\item The loops functor $\Omega\colon\Top_*\to\Top_*$ preserves products.
\end{enumerate}
\end{exercise}

\begin{exercise}\label{exer:adj-iff-representable}
\index{adjunction!and representability}
Show that $G\colon\mathcal{D}\to\mathcal{C}$ has a left adjoint
if and only if for each $A\in\mathcal{C}$, the functor
$\Hom_\mathcal{C}(A,G-)\colon\mathcal{D}\to\Set$ is representable.
\end{exercise}

\begin{exercise}\label{exer:reflective-subcategory}
\index{reflective subcategory}
\index{adjunction!reflective subcategory}
A full subcategory $\mathcal{D}\hookrightarrow\mathcal{C}$ is
\textbf{reflective} if the inclusion functor $\iota$ has a left adjoint
$L$ (the \textbf{reflector}).  Show that:
\begin{enumerate}[label=(\alph*)]
\item The counit $\varepsilon\colon L\iota\Rightarrow\Id_\mathcal{D}$
    is a natural isomorphism.
\item $\Ab\hookrightarrow\Grp$ is reflective (the reflector is abelianisation).
\item The full subcategory of complete metric spaces in $\mathsf{Met}$
    is reflective (the reflector is Cauchy completion).
\end{enumerate}
\end{exercise}

\begin{exercise}\label{exer:fixed-points-adjunction}
\index{adjunction!and fixed points}
Let $F\dashv G$ with unit~$\eta$ and counit~$\varepsilon$.
An object $A\in\mathcal{C}$ is called \textbf{$G$-split} if
$\eta_A$ is a split monomorphism.
\begin{enumerate}[label=(\alph*)]
\item Show that $\eta_A$ is a split mono for every~$A$ if and only if
    $G$ is faithful.
\item Show that $\eta$ is a natural isomorphism if and only if
    $G$ is fully faithful (i.e.\ $\mathcal{D}$ is a reflective
    subcategory of~$\mathcal{C}$ up to equivalence).
\end{enumerate}
\end{exercise}

\begin{exercise}\label{exer:GAFT-apply}
\index{General Adjoint Functor Theorem!exercises}
\begin{enumerate}[label=(\alph*)]
\item Use the GAFT to prove that the forgetful functor
    $U\colon\mathsf{CompHaus}\to\Top$ has a left adjoint
    (the Stone--\v{C}ech compactification~$\beta$).
    Verify the solution set condition.
\item Use the SAFT to show that any continuous (limit-preserving)
    functor $G\colon\Ab\to\Set$ has a left adjoint.
\end{enumerate}
\end{exercise}

\begin{exercise}\label{exer:adj-monad}
\index{monad!from adjunction}
\index{adjunction!associated monad}
Let $F\dashv G$ with unit $\eta$ and counit~$\varepsilon$.
Define $T = GF\colon\mathcal{C}\to\mathcal{C}$.
Show that $(T,\eta,\mu)$ is a \textbf{monad} on~$\mathcal{C}$,
where $\mu = G\varepsilon F\colon T^2 = GFGF\Rightarrow GF = T$.
That is, verify:
\begin{enumerate}[label=(\alph*)]
\item $\mu\circ T\mu = \mu\circ\mu T$ (associativity);
\item $\mu\circ T\eta = \mu\circ\eta T = \id_T$ (unit laws).
\end{enumerate}
\emph{Hint:} Use the triangle identities and naturality.
\[
\begin{tikzcd}
T^3 \arrow[r, "T\mu"] \arrow[d, "\mu T"'] & T^2 \arrow[d, "\mu"] \\
T^2 \arrow[r, "\mu"'] & T
\end{tikzcd}
\qquad
\begin{tikzcd}
T \arrow[r, "T\eta"] \arrow[dr, equal] & T^2 \arrow[d, "\mu"]
  & T \arrow[l, "\eta T"'] \arrow[dl, equal] \\
& T &
\end{tikzcd}
\]
(We will study monads in detail in a later chapter.)
\end{exercise}

\index{adjunction|)}
